\documentclass[../../../include/open-logic-section]{subfiles}
%READERNOTE{SOL6-A227: स्वसदस्यत्व-विरोधी युक्तिवाद प्रथम निर्माण-टप्पा आणि त्याच्या आधी तयार झालेले सदस्य वापरतो. सदस्यांचा शेवटचा टप्पा असावा किंवा पुढचा टप्पा लगतचा असावा अशी अतिरिक्त अट नाही.}

\begin{document}

\olfileid{sth}{z}{sep}
\olsection{विभक्तीकरण}

अनौपचारिक गुणधर्माधारित संचरचनेच्या जागी पुढील तत्त्व घेऊ:

\begin{axiom}[विभक्तीकरण-रूपबंध] प्रत्येक सूत्र $\phi(x)$ साठी पुढील एक स्वयंसिद्धक आहे: कोणताही $A$ दिल्यास $\Setabs{x \in A}{\phi(x)}$ हा संच अस्तित्वात असतो.
\end{axiom}

हे एकच स्वयंसिद्धक नाही, हे लक्षात घ्या. हा स्वयंसिद्धकांचा \emph{रूपबंध} आहे. विभक्तीकरणाची स्वयंसिद्धके \emph{अनंत} आहेत; प्रत्येक सूत्र $\phi(x)$ साठी एक. हा रूपबंध पुढीलप्रमाणेही लिहिता येतो (आणि सहसा तसाच लिहिला जातो):

\begin{defish}
``$S$'' नसलेल्या कोणत्याही सूत्र $\phi(x)$ साठी पुढील एक स्वयंसिद्धक आहे:
\[
	\forall A \exists S \forall x(x \in S \liff (\phi(x) \land x \in A)).
\]
\end{defish}

\olref[sth][][]{part} च्या सुरुवातीला सांगितलेल्या संकेतानुसार, विभक्तीकरणाच्या स्वयंसिद्धकांतील सूत्रांना~$\phi$ प्राचले असू शकतात.\footnote{याचा अर्थ काय, यासाठी \olref[sfr][infinite][induction]{natinductionschema} नंतरची चर्चा पाहा.}

आपण सांगत आलेली संचांची संचयी-क्रमिक संकल्पना विभक्तीकरणाचे लगेच समर्थन करते. ते पाहण्यासाठी $A$ हा संच मानू. \stageshier{} नुसार $A$ कोणत्या तरी टप्प्या~$S$ वर तयार होतो. $A$ टप्पा~$S$ वर तयार झाला असल्याने $A$ चे सर्व सदस्य टप्पा~$S$ च्या आधी तयार झालेले असतात (\stagesacc{} नुसार). आता $A$ चे सदस्य असलेले आणि $\phi$ पूर्ण करणारे सर्व संच घ्या; हेसुद्धा टप्पा~$S$ च्या आधीच तयार झालेले असतात. म्हणून \stagesacc{} नुसार टप्पा~$S$ वर त्यांचा $\Setabs{x \in A}{\phi(x)}$ हा संच तयार होतो.

अनौपचारिक गुणधर्माधारित संचरचनेच्या विपरीत, विभक्तीकरणामुळे रसेलची विरोधापत्ती टळते. कारण $\Setabs{x}{x \notin
x}$ हा संच आहे, असे थेट म्हणता येत नाही. पण एखादा संच~$A$ \emph{दिला} असेल, तर $R_A = \Setabs{x \in A}{x \notin x}$ हा संच आहे, असे म्हणता येते. यातून फक्त $R_A \notin R_A$ आणि $R_A \notin A$ हे निष्पन्न होते; त्यांत चिंतेचे काही नाही.

मात्र विभक्तीकरणाचा एक तात्काळ आणि महत्त्वाचा परिणाम आहे:

\begin{thm}\ollabel{thm:NoUniversalSet}
\emph{सर्वव्यापी} संच अस्तित्वात नाही; म्हणजे $\Setabs{x}{x = x}$ हा संच अस्तित्वात नाही.
\end{thm}

\begin{proof}
व्यतिरेकासाठी $V$ हा सर्वव्यापी संच आहे असे मानू. मग विभक्तीकरणानुसार $R =
\Setabs{x \in V}{x \notin x} = \Setabs{x}{x \notin x}$ हा संच अस्तित्वात असतो; पण याने रसेलची विरोधापत्ती निर्माण होते.
\end{proof}

सर्वव्यापी संच नसणे---खरे तर संचांच्या पदानुक्रमाला शेवट नसणे---ही संचांच्या संचयी-क्रमिक संकल्पनेमागील एक मूलभूत कल्पना आहे. त्यामुळे सहजज्ञानाने याच निष्कर्षापर्यंत दुसऱ्या मार्गाने कसे पोहोचता येते, हे पाहणे उपयुक्त ठरेल. सर्वव्यापी संच स्वतःचा !!a{element} असलाच पाहिजे. पण टप्पे सुक्रमित असल्याच्या गृहीतकानुसार प्रत्येक संच प्रथम तयार होण्याचा एक टप्पा असतो. त्या टप्प्यावर तयार होणाऱ्या संचाचे सर्व !!{element}s आधीच्या टप्प्यांवर तयार झालेले असतात. म्हणून \emph{कोणताही} संच स्वतःचा सदस्य नसतो: तो स्वतःचा सदस्य असता, तर प्रथम तयार होण्याच्या टप्प्याच्या आधीच तो तयार झाला असता; हा व्याघात आहे. या युक्तिवादाला सर्व सदस्यांचा एखादा शेवटचा टप्पा किंवा त्याचा लगतचा उत्तरवर्ती टप्पा आवश्यक नाही. (संचाच्या ``क्रमस्थानाची'' संकल्पना वापरून हा युक्तिवाद अधिक काटेकोर कसा करायचा, ते \olref[sth][spine][rank]{defnsetrank} मध्ये पाहू. मात्र त्यासाठी आधी आणखी काही स्वयंसिद्धके लागतील.)

विभक्तीकरण आणि विस्तारात्मकता यांचे आणखी काही परिणाम पाहू.

\begin{prop}\ollabel{prop:emptyexists}
कोणताही संच अस्तित्वात असेल, तर $\emptyset$ अस्तित्वात असतो.
\end{prop}

\begin{proof}
$A$ हा संच असेल, तर विभक्तीकरणानुसार $\emptyset = \Setabs{x \in A}{x \neq x}$ अस्तित्वात असतो.
\end{proof}

\begin{prop}
कोणतेही संच $A$ आणि $B$ यांच्यासाठी $A \setminus B$ अस्तित्वात असतो.
\end{prop}

\begin{proof}
विभक्तीकरणानुसार $A \setminus B = \Setabs{x \in A}{x \notin B}$ अस्तित्वात असतो.
\end{proof}

याशिवाय (जवळजवळ) मनमानी छेदही अस्तित्वात असतात:

\begin{prop}\ollabel{prop:intersectionsexist}
$A \neq \emptyset$ असेल, तर $\bigcap A = \Setabs{x}{(\forall y \in A)x \in y}$ अस्तित्वात असतो.
\end{prop}

\begin{proof}
$A \neq \emptyset$ मानू; मग एखादा $c \in A$ असतो. म्हणून $\bigcap A =
\Setabs{x}{(\forall y \in A)x \in y} = \Setabs{x \in c}{(\forall y \in
A)x \in y}$; हा संच विभक्तीकरणानुसार अस्तित्वात असतो.
\end{proof}

पण $A \neq \emptyset$ ही अट लक्षात घ्या. कारण $\bigcap
\emptyset$ हा रिक्ततेमुळे सर्वव्यापी संच ठरेल; ते \olref{thm:NoUniversalSet} शी विसंगत आहे.

\end{document}
